\section{MARL Agent}

\subsection{Overview}

The \texttt{marl/} directory implements a multi-agent reinforcement learning approach using graph observations, graph convolutional networks, self-attention, and PPO. The intent is to represent the cyber environment as a changing network rather than as a fixed flat vector.

\subsection{Graph Observation Wrapper}

\texttt{marl/wrappers/graph\_wrapper.py} extends \texttt{EnterpriseMAE}. It translates CybORG observations into PyTorch-Geometric-compatible graph states and translates learned action IDs back into CybORG actions.

The graph contains typed nodes such as:
\begin{itemize}
  \item systems: hosts, servers, routers;
  \item connections: processes and open network connections;
  \item files: observed files, including suspicious malware-like artifacts;
  \item internet or structural nodes.
\end{itemize}

The wrapper also fuses graph features with tabular features from the blue wrapper, including mission phase, subnet policy, block state, and message-derived features. For explainability, it adds metadata such as chosen action type, chosen action string, and SHAP-oriented feature summaries.

\subsection{Action Space}

The learned policy emits one categorical distribution over a concatenated action vector:

\begin{itemize}
  \item node actions: \texttt{Analyse}, \texttt{Remove}, \texttt{Restore}, \texttt{DeployDecoy};
  \item edge actions: \texttt{AllowTraffic}, \texttt{BlockTraffic};
  \item global action: \texttt{Monitor}.
\end{itemize}

\texttt{GraphWrapper.action\_translator} maps the selected action ID to a CybORG action object for the correct subnet and host.

\subsection{Actor-Critic Architecture}

\texttt{marl/models/cage4.py} defines:

\begin{itemize}
  \item \texttt{InductiveActorNetwork};
  \item \texttt{InductiveCriticNetwork};
  \item \texttt{InductiveGraphPPOAgent};
  \item memory and batching utilities.
\end{itemize}

The actor applies two GCN layers. At each stage, a self-attention mechanism summarizes host and router embeddings into a global vector. The mission phase is embedded and injected into this global state. Action heads then score node, edge, and global actions.

\begin{lstlisting}
Graph state + phase
  -> GCN layers
  -> self-attention global summary
  -> node-action head
  -> edge-action head
  -> global-action head
  -> categorical policy
\end{lstlisting}

The critic uses a similar graph-processing structure to estimate value for PPO updates.

\subsection{Training and Evaluation}

Typical commands from the local workflow are:

\begin{lstlisting}
python marl/train.py mini_curriculum_smoke --mini \
  --strategy curriculum --warmup_updates 1 \
  --log_per_profile --device cpu
\end{lstlisting}

For GPU experiments:

\begin{lstlisting}
python marl/train.py exp_curriculum_gpu \
  --KIServer --strategy curriculum \
  --warmup_updates 200 --harden_updates 200 \
  --log_per_profile --phase_reward_mode red_only \
  --device cuda
\end{lstlisting}

The profile registry is tolerant: profiles missing from the checked-out tree are ignored and remaining weights are renormalized.
